\section{Cohomología de los esquemas afines}
\label{section:III.1}
\label{section:III.1-fr}

\subsection{Repaso del complejo del álgebra exterior}
\label{subsection:III.1.1}
\label{subsection:III.1.1-fr}

\begin{env}[1.1.1]
\label{III.1.1.1}
\label{III.1.1.1-fr}
Sea $A$ un anillo y sea $\mathbf f=(f_i)_{1\leq i\leq r}$ un sistema de $r$ elementos de $A$.
El \emph{complejo del álgebra exterior $K_\bullet(\mathbf f)$} correspondiente a $\mathbf f$ es un complejo de cadenas (G, I, 2.2) definido del modo siguiente: el $A$-módulo graduado $K_\bullet(\mathbf f)$ es igual al \emph{álgebra exterior $\bigwedge(A^r)$}, graduada de la manera usual, y la aplicación borde es la \emph{multiplicación interior $i_{\mathbf f}$} por $\mathbf f$ considerado como elemento del dual $(A^r)\check{}$; se recuerda que $i_{\mathbf f}$ es una \emph{antiderivación} de grado $-1$ de $\bigwedge(A^r)$, y que, si $(\mathbf e_i)_{1\leq i\leq r}$ es la base canónica de $A^r$, se tiene $i_{\mathbf f}(\mathbf e_i)=f_i$; la verificación de la condición $i_{\mathbf f}\circ i_{\mathbf f}=0$ es inmediata.

Una definición equivalente es la siguiente: para cada $i$, se considera un complejo de cadenas $K_\bullet(f_i)$ definido por $K_0(f_i)=K_1(f_i)=A$, $K_n(f_i)=0$ para $n\neq0,1$; la aplicación borde queda definida por la condición de que $d_1:A\to A$ sea la \emph{multiplicación por $f_i$}.
Se toma entonces como $K_\bullet(\mathbf f)$ el \emph{producto tensorial $K_\bullet(f_1)\otimes K_\bullet(f_2)\otimes\ldots\otimes K_\bullet(f_r)$} (G, I, 2.7), provisto de su grado total; la verificación del isomorfismo de este complejo con el definido anteriormente es inmediata.
\end{env}

\begin{env}[1.1.2]
\label{III.1.1.2}
\label{III.1.1.2-fr}
Para todo $A$-módulo $M$, se define el \emph{complejo de cadenas}
\begin{equation}
\label{III.1.1.2.1-fr}
\tag{1.1.2.1}
K_\bullet(\mathbf f,M)=K_\bullet(\mathbf f)\otimes_A M
\end{equation}
y el \emph{complejo de cocadenas} (G, I, 2.2)
\begin{equation}
\label{III.1.1.2.2-fr}
\tag{1.1.2.2}
K^\bullet(\mathbf f,M)=\operatorname{Hom}_A(K_\bullet(\mathbf f),M).
\end{equation}

Si $g$ es una $k$-cocadena de este último complejo y se pone
\[
g(i_1,\ldots,i_k)=g(\mathbf e_{i_1}\wedge\ldots\wedge\mathbf e_{i_k}),
\]
entonces $g$ se identifica con una aplicación \emph{alternada} de $[1,r]^k$ en $M$, y de las definiciones anteriores se deduce que
\begin{equation}
\label{III.1.1.2.3-fr}
\tag{1.1.2.3}
d^kg(i_1,i_2,\ldots,i_{k+1})=
\sum_{h=1}^{k+1}(-1)^{h-1}f_{i_h}
g(i_1,\ldots,\widehat{i}_h,\ldots,i_{k+1}).
\end{equation}
\end{env}

\oldpage[III]{83}

\begin{env}[1.1.3]
\label{III.1.1.3-fr}
De los complejos precedentes se deducen, como de costumbre, los \emph{$A$-módulos
de homología y de cohomología} (G, I, 2.2)
\begin{align}
\label{III.1.1.3.1-fr}
\tag{1.1.3.1}
H_\bullet(\mathbf f,M)&=H_\bullet(K_\bullet(\mathbf f,M))\\
\label{III.1.1.3.2-fr}
\tag{1.1.3.2}
H^\bullet(\mathbf f,M)&=H^\bullet(K^\bullet(\mathbf f,M)).
\end{align}
Se define además un $A$-\emph{isomorfismo}
$K_\bullet(\mathbf f,M)\xrightarrow{\sim}K^\bullet(\mathbf f,M)$ haciendo
corresponder a toda cadena
$z=\sum(\mathbf e_{i_1}\wedge\ldots\wedge\mathbf e_{i_k})\otimes
z_{i_1\ldots i_k}$ la cocadena $g_z$ tal que
$g_z(j_1,\ldots,j_{r-k})=\varepsilon z_{i_1\ldots i_k}$, donde
$(j_h)_{1\leq h\leq r-k}$ es la sucesión estrictamente creciente complementaria
de la sucesión estrictamente creciente $(i_h)_{1\leq h\leq k}$ en $[1,r]$ y
$\varepsilon=(-1)^\nu$, siendo $\nu$ el número de inversiones de la permutación
$i_1,\ldots,i_k,j_1,\ldots,j_{r-k}$ de $[1,r]$. Se verifica que
$g_{dz}=d(g_z)$, lo que da un isomorfismo
\begin{equation}
\label{III.1.1.3.3-fr}
\tag{1.1.3.3}
H^i(\mathbf f,M)\xrightarrow{\sim}H_{r-i}(\mathbf f,M)
\qquad\text{para }0\leq i\leq r.
\end{equation}
En este capítulo tendremos que considerar sobre todo los módulos de cohomología
$H^\bullet(\mathbf f,M)$.

Para un $\mathbf f$ dado, es inmediato (G, I, 2.1) que
$M\leadsto H^\bullet(\mathbf f,M)$ es un \emph{funtor cohomológico}
(T, II, 2.1), de la categoría de los $A$-módulos en la de los $A$-módulos
graduados, nulo para los grados $<0$ y $>r$. Además, se tiene
\begin{equation}
\label{III.1.1.3.4-fr}
\tag{1.1.3.4}
H^0(\mathbf f,M)=\operatorname{Hom}_A(A/(\mathbf f),M)
\end{equation}
designando por $(\mathbf f)$ el ideal de $A$ engendrado por
$f_1,\ldots,f_r$; esto resulta inmediatamente de
(\hyperref[III.1.1.2.3-fr]{1.1.2.3}), y está claro que
$H^0(\mathbf f,M)$ se identifica con el submódulo de $M$ \emph{anulado por
$(\mathbf f)$}. Análogamente, se tiene según
(\hyperref[III.1.1.2.3-fr]{1.1.2.3})
\begin{equation}
\label{III.1.1.3.5-fr}
\tag{1.1.3.5}
H^r(\mathbf f,M)=M/\Bigl(\sum_{i=1}^r f_iM\Bigr)
=(A/(\mathbf f))\otimes_A M.
\end{equation}
Utilizaremos el siguiente resultado conocido, del que recordaremos una
demostración para ser completos:
\end{env}

\begin{proposition}[1.1.4]
\label{III.1.1.4}
\label{III.1.1.4-fr}
Sean $A$ un anillo, $\mathbf f=(f_i)_{1\leq i\leq r}$ una familia finita de elementos de $A$ y $M$ un $A$-módulo.
Si, para $1\leq i\leq r$, la aplicación $z\leadsto f_i.z$ sobre $M_{i-1}=M/(f_1M+\ldots+f_{i-1}M)$ es inyectiva, entonces $H^i(\mathbf f,M)=0$ para $i\neq r$.
\end{proposition}

\begin{proof}
Basta probar que $H_i(\mathbf f,M)=0$ para todo $i>0$, en virtud de (\hyperref[III.1.1.3.3]{1.1.3.3}).
Razonemos por inducción sobre $r$, siendo trivial el caso $r=0$ siendo trivial. Pongamos $\mathbf f'=(f_i)_{1\leq i\leq r-1}$; esta familia satisface las condiciones del enunciado, de modo que, si se pone $L_\bullet=K_\bullet(\mathbf f',M)$, por hipótesis se tiene $H_i(L_\bullet)=0$ para $i>0$, y $H_0(L_\bullet)=M_{r-1}$ en virtud de (\hyperref[III.1.1.3.3]{1.1.3.3}) y (\hyperref[III.1.1.3.5]{1.1.3.5}).
Para abreviar, pongamos $K_\bullet=K_\bullet(f_r)=K_0\oplus K_1$, con $K_0=K_1=A$ y $d_1:K_1\to K_0$ igual a la multiplicación por $f_r$; por definición \sref{III.1.1.1}, se tiene $K_\bullet(\mathbf f,M)=K_\bullet\otimes_A L_\bullet$.
Se dispone del siguiente lema:

\begin{lemma}[1.1.4.1]
\label{III.1.1.4.1}
\label{III.1.1.4.1-fr}
Sea $K_\bullet$ un complejo de cadenas de $A$-módulos libres, nulo salvo en las dimensiones $0$ y $1$.
Para todo complejo de cadenas $L_\bullet$ de $A$-módulos se tiene una sucesión exacta
\[
0\longrightarrow H_0(K_\bullet\otimes H_p(L_\bullet))
\longrightarrow H_p(K_\bullet\otimes L_\bullet)
\longrightarrow H_1(K_\bullet\otimes H_{p-1}(L_\bullet))
\longrightarrow0
\]
para todo índice $p$.
\end{lemma}

\oldpage[III]{84}
Este es un caso particular de una sucesión exacta de términos de grado bajo de la sucesión espectral de K\"unneth (M, XVII, 5.2 \emph{(a)} y G, I, 5.5.2); puede demostrarse directamente como sigue.
Consideremos $K_0$ y $K_1$ como complejos de cadenas (nulos en las dimensiones $\neq0$ y $\neq1$, respectivamente); se tiene entonces una sucesión exacta de complejos
\[
0\longrightarrow K_0\otimes L_\bullet
\longrightarrow K_\bullet\otimes L_\bullet
\longrightarrow K_1\otimes L_\bullet\longrightarrow0
\] a la cual podemos aplicar la sucesión exacta de homología \[
\ldots\longrightarrow H_{p+1}(K_1\otimes L_\bullet)
\xrightarrow{\partial}H_p(K_0\otimes L_\bullet)
\longrightarrow H_p(K_\bullet\otimes L_\bullet)
\longrightarrow H_p(K_1\otimes L_\bullet)
\xrightarrow{\partial}H_{p-1}(K_0\otimes L_\bullet)
\longrightarrow\ldots
\]
Pero es evidente que $H_p(K_0\otimes L_\bullet)=K_0\otimes H_p(L_\bullet)$ y $H_p(K_1\otimes L_\bullet)=K_1\otimes H_{v-1}(L_\bullet)$ para todo $p$; además, se comprueba inmediatamente que el operador $\partial:K_1\otimes H_p(L_\bullet)\to K_0\otimes H_p(L_\bullet)$ no es otro que $d_1\otimes1$; el lema se deduce así de la sucesión exacta anterior y de la definición de $H_0(K\otimes H_p(L_\bullet))$ y $H_1(K_\bullet\otimes H_{p-1}(L_\bullet))$.

Establecido el lema, el final de la demostración de la Proposición \sref{III.1.1.4} es inmediato: la hipótesis de inducción y el Lema \sref{III.1.1.4.1} dan $H_p(K_\bullet\otimes L_\bullet)=0$ para $p\geq2$; además, si se prueba que $H_1(K_\bullet,H_0(L_\bullet))=0$, se deduce también del Lema \sref{III.1.1.4.1} que $H_1(K_\bullet\otimes L_\bullet)=0$; pero, por definición, $H_1(K_\bullet,H_0(L_\bullet))$ no es otro que el núcleo de la aplicación $z\leadsto f_r.z$ sobre $M_{r-1}$, y, como por hipótesis este núcleo es nulo, queda terminada la demostración.
\end{proof}

\begin{env}[1.1.5]
\label{III.1.1.5}
\label{III.1.1.5-fr}
Sea $\mathbf g=(g_i)_{1\leq i\leq r}$ una segunda sucesión de $r$ elementos de $A$, y pongamos $\mathbf f\mathbf g=(f_ig_i)_{1\leq i\leq r}$.
Puede definirse un homomorfismo canónico de complejos
\begin{equation}
\label{III.1.1.5.1-fr}
\tag{1.1.5.1}
\varphi_{\mathbf g}:K_\bullet(\mathbf f\mathbf g)\longrightarrow
K_\bullet(\mathbf f)
\end{equation}
como la extensión canónica al álgebra exterior $\bigwedge(A^r)$ de la aplicación $A$-lineal $(x_1,\ldots,x_r)\leadsto(g_1x_1,\ldots,g_rx_r)$ de $A^r$ en sí mismo.
Para ver que se trata de un homomorfismo de complejos, basta observar, en general, que, si $u:E\to F$ es una aplicación $A$-lineal, $\mathbf x\in\check F$ y $\mathbf y={}^tu(\mathbf x)\in\check E$, entonces se tiene la fórmula
\begin{equation}
\label{III.1.1.5.2-fr}
\tag{1.1.5.2}
(\bigwedge u)\circ i_{\mathbf y}=i_{\mathbf x}\circ(\bigwedge u)~;
\end{equation}
en efecto, los dos miembros son antiderivaciones de $\bigwedge F$, y basta verificar que coinciden en $F$, lo cual se deduce inmediatamente de las definiciones.

Cuando se identifica $K_\bullet(\mathbf f)$ con el producto tensorial de los $K_\bullet(f_i)$ \sref{III.1.1.1}, $\varphi_{\mathbf g}$ es el producto tensorial de los $\varphi_{g_i}$, donde $\varphi_{g_i}$ es la identidad en grado $0$ y la multiplicación por $g_i$ en grado $1$.
\end{env}

\begin{env}[1.1.6]
\label{III.1.1.6-fr}
En particular, para todo par de enteros $m,n$ tales que $0\leq n\leq m$, se tienen homomorfismos de complejos
\begin{equation}
\label{III.1.1.6.1-fr}
\tag{1.1.6.1}
\varphi_{\mathbf f^{m-n}}:K_\bullet(\mathbf f^m)\longrightarrow
K_\bullet(\mathbf f^n)
\end{equation}
y por consiguiente homomorfismos
\begin{align}
\label{III.1.1.6.2-fr}
\tag{1.1.6.2}
\varphi_{\mathbf f^{m-n}} &:K^\bullet(\mathbf f^n,M)\longrightarrow
K^\bullet(\mathbf f^m,M)\\
\label{III.1.1.6.3-fr}
\tag{1.1.6.3}
\varphi_{\mathbf f^{m-n}} &:H^\bullet(\mathbf f^n,M)\longrightarrow
H^\bullet(\mathbf f^m,M).
\end{align}

\oldpage[III]{85}
Estos últimos satisfacen evidentemente la condición de transitividad
$\varphi_{\mathbf f^{m-p}}=
\varphi_{\mathbf f^{m-n}}\circ\varphi_{\mathbf f^{n-p}}$ para $p\leq n\leq m$; definen, por tanto, dos \emph{sistemas inductivos} de $A$-módulos; se pondrá
\begin{align}
\label{III.1.1.6.4-fr}
\tag{1.1.6.4}
C^\bullet((\mathbf f),M)&=\varinjlim_n K^\bullet(\mathbf f^n,M)\\
\label{III.1.1.6.5-fr}
\tag{1.1.6.5}
H^\bullet((\mathbf f),M)&=H^\bullet(C^\bullet((\mathbf f),M))
=\varinjlim_n H^\bullet(\mathbf f^n,M)
\end{align}
proviniendo la última igualdad de que el paso al límite inductivo conmuta con el funtor $H^\bullet$ (G, I, 2.1).
Se verá más adelante \sref{III.1.4.3} que $H^\bullet((\mathbf f),M)$ depende en realidad solo del \emph{ideal $(\mathbf f)$} de $A$ (e incluso solo de la topología $(\mathbf f)$-preádica sobre $A$), lo que justifica las notaciones.

Es claro que $M\leadsto C^\bullet((\mathbf f),M)$ es un funtor exacto $A$-lineal y que $M\leadsto H^\bullet((\mathbf f),M)$ es un funtor cohomológico.
\end{env}

\begin{env}[1.1.7]
\label{III.1.1.7}
\label{III.1.1.7-fr}
Pónganse $\mathbf f=(f_i)\in A^r$ y $\mathbf g=(g_i)\in A^r$; designemos por $e_{\mathbf g}$ la multiplicación por la izquierda por el vector $\mathbf g\in A^r$ sobre el álgebra exterior $\bigwedge(A^r)$; se sabe que se tiene la \emph{fórmula de homotopía}
\begin{equation}
\label{III.1.1.7.1-fr}
\tag{1.1.7.1}
i_{\mathbf f}e_{\mathbf g}+e_{\mathbf g}i_{\mathbf f}
=\langle\mathbf g,\mathbf f\rangle 1
\end{equation}
en el $A$-módulo $A^r$ ($1$ denota el automorfismo identidad de $A^r$); esta relación implica también que \emph{en el complejo $K_\bullet(\mathbf f)$} se tiene
\begin{equation}
\label{III.1.1.7.2-fr}
\tag{1.1.7.2}
de_{\mathbf g}+e_{\mathbf g}d=\langle\mathbf g,\mathbf f\rangle 1.
\end{equation}

Si el ideal $(\mathbf f)$ es igual a $A$, existe $\mathbf g\in A^r$ tal que $\langle\mathbf g,\mathbf f\rangle=\sum_{i=1}^r g_if_i=1$.
Por consiguiente (G, I, 2.4):
\end{env}

\begin{proposition}[1.1.8]
\label{III.1.1.8}
\label{III.1.1.8-fr}
Supongamos que el ideal $(\mathbf f)$ generado por los $f_i$ es igual a $A$.
Entonces el complejo $K_\bullet(\mathbf f)$ es homotópicamente trivial, al igual que los complejos $K_\bullet(\mathbf f,M)$ y $K^\bullet(\mathbf f,M)$ para todo $A$-módulo $M$.
\end{proposition}

\begin{corollary}[1.1.9]
\label{III.1.1.9}
\label{III.1.1.9-fr}
Si $(\mathbf f)=A$, entonces $H^\bullet(\mathbf f,M)=0$ y $H^\bullet((\mathbf f),M)=0$ para todo $A$-módulo $M$.
\end{corollary}

\begin{proof}
En efecto, se tiene entonces $(\mathbf f^n)=A$ para todo $n$.
\end{proof}

\begin{remark}[1.1.10]
\label{III.1.1.10}
\label{III.1.1.10-fr}
Con las mismas notaciones anteriores, sean $X=\operatorname{Spec}(A)$ y sea $Y$ el subesquema cerrado de $X$ definido por el ideal $(\mathbf f)$. Probaremos en el \S~9 que $H^\bullet((\mathbf f),M)$ es isomorfo a la cohomología $H_Y^\bullet(X,\widetilde M)$ correspondiente al antifiltro $\Phi$ de subconjuntos cerrados de $Y$ (T, 3.2). Mostraremos también que la prop. (\hyperref[III.1.2.3-fr]{1.2.3}) aplicada a $X$ y a $\mathcal F=\widetilde M$, es un caso particular de una sucesión exacta de cohomología
\[
\ldots\longrightarrow H_Y^p(X,\mathcal F)\longrightarrow H^p(X,\mathcal F)
\longrightarrow H^p(X\setmin Y,\mathcal F)\longrightarrow H_Y^{p+1}(X,\mathcal F)
\longrightarrow\ldots
\]
\end{remark}

\subsection{Cohomología de \v Cech de un recubrimiento abierto}
\label{subsection:III.1.2}
\label{subsection:III.1.2-fr}

\begin{notation}[1.2.1]
\label{III.1.2.1}
\label{III.1.2.1-fr}
En esta sección se denota por:
\begin{itemize}
  \item $X$ un preesquema;
  \item $\mathcal F$ un $\mathcal O_X$-módulo cuasicoherente;
\oldpage[III]{86}
  \item $A=\Gamma(X,\mathcal O_X)$, $M=\Gamma(X,\mathcal F)$;
  \item $\mathbf f=(f_i)_{1\leq i\leq r}$ un sistema finito de elementos de $A$;
  \item $U_i=X_{f_i}$, el abierto \sref[0]{0.5.5.2} de los $x\in X$ tales que $f_i(x)\neq0$;
  \item $U=\bigcup_{i=1}^r U_i$;
  \item $\mathfrak U$ el recubrimiento $(U_i)_{1\leq i\leq r}$ de $U$.
\end{itemize}
\end{notation}

\begin{env}[1.2.2]
\label{III.1.2.2}
\label{III.1.2.2-fr}
Supongamos que $X$ es un preesquema cuyo espacio subyacente es \emph{noetheriano}, o bien un \emph{esquema} cuyo espacio subyacente es \emph{cuasicompacto}.
Se sabe entonces \sref[I]{I.9.3.3} que $\Gamma(U_i,\mathcal F)=M_{f_i}$. Pondremos \[
U_{i_0i_1\ldots i_p}=\bigcap_{k=0}^p U_{i_k}
=X_{f_{i_0}f_{i_1}\ldots f_{i_p}}
\]
\sref[0]{0.5.5.3}; se tiene, por tanto, también
\begin{equation}
\label{III.1.2.2.1-fr}
\tag{1.2.2.1}
\Gamma(U_{i_0i_1\ldots i_p},\mathcal F)=M_{f_{i_0}f_{i_1}\ldots f_{i_p}}
\end{equation}

Por \sref[0]{0.1.6.1}, $M_{f_{i_0}f_{i_1}\ldots f_{i_p}}$ se identifica con el límite inductivo $\varinjlim_n M_{i_0i_1\ldots i_p}^{(n)}$, donde el sistema inductivo está formado por los $M_{i_0i_1\ldots i_p}^{(n)}=M$, siendo los homomorfismos $\varphi_{nm}:M_{i_0i_1\ldots i_p}^{(m)}\to
M_{i_0i_1\ldots i_p}^{(n)}$ la multiplicación por $(f_{i_0}f_{i_1}\ldots f_{i_p})^{n-m}$ para $m\leq n$. Designemos por $C_n^p(M)$ el conjunto de aplicaciones \emph{alternadas} de $[1,r]^{p+1}$ en $M$ (para todo $n$); estos $A$-módulos forman también un sistema inductivo respecto de los $\varphi_{nm}$.
Si $C^p(\mathfrak U,\mathcal F)$ es el grupo de las $p$-cocadenas \emph{alternadas} de \v Cech relativas al recubrimiento $\mathfrak U$, con coeficientes en $\mathcal F$ (G, II, 5.1), de lo anterior se deduce que puede escribirse
\begin{equation}
\label{III.1.2.2.2-fr}
\tag{1.2.2.2}
C^p(\mathfrak U,\mathcal F)=\varinjlim_n C_n^p(M).
\end{equation}

Con las notaciones de \sref{III.1.1.2}, $C_n^p(M)$ se identifica con $K^{p+1}(\mathbf f^n,M)$, y la aplicación $\varphi_{nm}$ se identifica con la aplicación $\varphi_{\mathbf f^{n-m}}$ definida en \sref{III.1.1.6}.
Se tiene así, para todo $p\geq0$, un isomorfismo canónico funtorial en $\mathcal F$
\begin{equation}
\label{III.1.2.2.3-fr}
\tag{1.2.2.3}
C^p(\mathfrak U,\mathcal F)\xrightarrow{\sim}C^{p+1}((\mathbf f),M).
\end{equation}

Además, la fórmula (\hyperref[III.1.1.2.3]{1.1.2.3}) y la definición de la cohomología de un recubrimiento (G, II, 5.1) muestran que los isomorfismos (\hyperref[III.1.2.2.3]{1.2.2.3}) son compatibles con las aplicaciones coborde.
\end{env}

\begin{proposition}[1.2.3]
\label{III.1.2.3}
\label{III.1.2.3-fr}
Si $X$ es un preesquema cuyo espacio subyacente es noetheriano, o un esquema cuyo espacio subyacente es cuasicompacto, existe un isomorfismo canónico funtorial en $\mathcal F$
\begin{equation}
\label{III.1.2.3.1-fr}
\tag{1.2.3.1}
H^p(\mathfrak U,\mathcal F)\xrightarrow{\sim}H^{p+1}((\mathbf f),M)
\qquad\text{para todo }p\geq1.
\end{equation}

Además, se tiene una sucesión exacta funtorial en $\mathcal F$
\begin{equation}
\label{III.1.2.3.2-fr}
\tag{1.2.3.2}
0\longrightarrow H^0((\mathbf f),M)\longrightarrow M
\longrightarrow H^0(\mathfrak U,\mathcal F)
\longrightarrow H^1((\mathbf f),M)\longrightarrow0.
\end{equation}
\end{proposition}

\oldpage[III]{87}
\begin{proof}
Los isomorfismos (\hyperref[III.1.2.3.1]{1.2.3.1}) son consecuencias inmediatas de lo visto en \sref{III.1.2.2}.
Por otra parte, se tiene $C^0(\mathfrak U,\mathcal F)=C^1((\mathbf f),M)$; por consiguiente, $H^0(\mathfrak U,\mathcal F)$ se identifica con el subgrupo de los $1$-cociclos de $C^1((\mathbf f),M)$; como $M=C^0((\mathbf f),M)$, la sucesión exacta (\hyperref[III.1.2.3.2]{1.2.3.2}) no es otra que la dada por la definición de los grupos de cohomología $H^0((\mathbf f),M)$ y $H^1((\mathbf f),M)$.
\end{proof}

\begin{corollary}[1.2.4]
\label{III.1.2.4}
\label{III.1.2.4-fr}
Supongamos que los $X_{f_i}$ son cuasicompactos y que existen $g_i\in\Gamma(U,\mathcal F)$ tales que
$\sum_i g_i(f_i\vert U)=1\vert U$. Entonces, para todo
$(\mathcal O_X\vert U)$-Módulo cuasicoherente $\mathcal G$, se tiene $H^p(\mathfrak U,\mathcal G)=0$ para $p>0$; si además $U=X$, el homomorfismo canónico \textup{(\hyperref[III.1.2.3.2]{1.2.3.2})} $M\to H^0(\mathfrak U,\mathcal F)$ es biyectivo.
\end{corollary}

\begin{proof}
Como, por hipótesis, los $U_i=X_{f_i}$ son cuasicompactos, también lo es $U$, y puede reducirse al caso $U=X$; la hipótesis implica entonces
$H^p((\mathbf f),M)=0$ para todo $p\geq0$ \sref{III.1.1.9}.
El corolario se deduce inmediatamente de (\hyperref[III.1.2.3.1]{1.2.3.1}) y (\hyperref[III.1.2.3.2]{1.2.3.2}).
\end{proof}

Se observará que, puesto que $H^0(\mathfrak U,\mathcal F)=H^0(U,\mathcal F)$ (G, II, 5.2.2), se ha demostrado de nuevo \sref[I]{I.1.3.7} como caso particular.

\begin{remark}[1.2.5]
\label{III.1.2.5}
\label{III.1.2.5-fr}
Supongamos que $X$ es un \emph{esquema afín}; entonces los $U_i=X_{f_i}=D(f_i)$ son abiertos afines, al igual que los $U_{i_0i_1\ldots i_p}$ (pero $U$ no es necesariamente afín).
En este caso, los funtores $\Gamma(X,\mathcal F)$ y $\Gamma(U_{i_0i_1\ldots i_p},\mathcal F)$ son exactos en $\mathcal F$ \sref[I]{I.1.3.11}.
Si se tiene una sucesión exacta $0\to\mathcal F'\to\mathcal F\to\mathcal F''\to0$ de $\mathcal O_X$-módulos cuasicoherentes, la sucesión de complejos
\[
0\longrightarrow C^\bullet(\mathfrak U,\mathcal F')
\longrightarrow C^\bullet(\mathfrak U,\mathcal F)
\longrightarrow C^\bullet(\mathfrak U,\mathcal F'')
\longrightarrow0
\]
es exacta y da, por tanto, una sucesión exacta de cohomología
\[
\ldots\longrightarrow H^p(\mathfrak U,\mathcal F')
\longrightarrow H^p(\mathfrak U,\mathcal F)
\longrightarrow H^p(\mathfrak U,\mathcal F'')
\xrightarrow{\partial}H^{p+1}(\mathfrak U,\mathcal F')
\longrightarrow\ldots
\]
Por otra parte, si se pone $M'=\Gamma(X,\mathcal F')$ y $M''=\Gamma(X,\mathcal F'')$, la sucesión $0\to M'\to M\to M''\to0$ es exacta; como $C^\bullet((\mathbf f),M)$ es un funtor exacto en $M$, se tiene asimismo la sucesión exacta de cohomología
\[
\ldots\longrightarrow H^p((\mathbf f),M')
\longrightarrow H^p((\mathbf f),M)
\longrightarrow H^p((\mathbf f),M'')
\xrightarrow{\partial}H^{p+1}((\mathbf f),M')
\longrightarrow\ldots
\]

Así pues, como el diagrama
\[
\xymatrix{
0\ar[r]&C^\bullet(\mathfrak U,\mathcal F')\ar[r]\ar[d]
&C^\bullet(\mathfrak U,\mathcal F)\ar[r]\ar[d]
&C^\bullet(\mathfrak U,\mathcal F'')\ar[r]\ar[d]&0\\
0\ar[r]&C^\bullet((\mathbf f),M')\ar[r]
&C^\bullet((\mathbf f),M)\ar[r]
&C^\bullet((\mathbf f),M'')\ar[r]&0
}
\]
\oldpage[III]{88}
es conmutativo, se concluye que los diagramas
\begin{equation}
\label{III.1.2.5.1-fr}
\tag{1.2.5.1}
\xymatrix{
H^p(\mathfrak U,\mathcal F'')\ar[r]^{\partial}\ar[d]
&H^{p+1}(\mathfrak U,\mathcal F')\ar[d]\\
H^{p+1}((\mathbf f),M'')\ar[r]_{\partial}
&H^{p+2}((\mathbf f),M')
}
\end{equation}
son conmutativos para todo $p$~(G, I, 2.1.1).
\end{remark}

\subsection{Cohomología de un esquema afín}
\label{subsection:III.1.3}
\label{subsection:III.1.3-fr}

\begin{theorem}[1.3.1]
\label{III.1.3.1}
\label{III.1.3.1-fr}
Sea $X$ un esquema afín.
Para todo $\mathcal O_X$-módulo cuasicoherente $\mathcal F$, se tiene $H^p(X,\mathcal F)=0$ para todo $p>0$.
\end{theorem}

\begin{proof}
Sea $\mathfrak U$ un recubrimiento finito de $X$ por los abiertos afines $X_{f_i}=D(f_i)$ ($1\leq i\leq r$); se sabe entonces que el ideal de $A=\Gamma(X,\mathcal O_X)$ generado por los $f_i$ es igual a $A$.
Del Corolario~\sref{III.1.2.4} se concluye, pues, que $H^p(\mathfrak U,\mathcal F)=0$ para $p>0$.
Como existen recubrimientos finitos de $X$ por abiertos afines arbitrariamente finos \sref[I]{I.1.1.10}, la definición de la cohomología de \v Cech (G, II, 5.8) muestra que también $\check H^p(X,\mathcal F)=0$ para $p>0$.
Pero esto se aplica asimismo a todo preesquema $X_f$, para $f\in A$ \sref[I]{I.1.3.6}; por tanto, $\check H^p(X_f,\mathcal F)=0$ para $p>0$.
Como $X_f\cap X_g=X_{fg}$, se deduce también que $H^p(X,\mathcal F)=0$ para todo $p>0$, en virtud de (G, II, 5.9.2).
\end{proof}

\begin{corollary}[1.3.2]
\label{III.1.3.2}
\label{III.1.3.2-fr}
Sea $Y$ un preesquema y sea $f:X\to Y$ un morfismo afín \sref[II]{II.1.6.1}.
Para todo $\mathcal O_X$-módulo cuasicoherente $\mathcal F$, se tiene $R^qf_*(\mathcal F)=0$ para $q>0$.
\end{corollary}

\begin{proof}
Por definición, $R^qf_*(\mathcal F)$ es el $\mathcal O_Y$-módulo asociado al prehaz $U\leadsto H^q(f^{-1}(U),\mathcal F)$, donde $U$ recorre los abiertos de $Y$.
Pero los abiertos afines forman una base de $Y$, y, para un abierto tal $U$, $f^{-1}(U)$ es afín \sref[II]{II.1.3.2}; por tanto, $H^q(f^{-1}(U),\mathcal F)=0$ por el Teorema~\sref{III.1.3.1}, lo que prueba el corolario.
\end{proof}

\begin{corollary}[1.3.3]
\label{III.1.3.3}
\label{III.1.3.3-fr}
Sea $Y$ un preesquema y sea $f:X\to Y$ un morfismo afín.
Para todo $\mathcal O_X$-módulo cuasicoherente $\mathcal F$, el homomorfismo canónico $H^p(Y,f_*(\mathcal F))\to H^p(X,\mathcal F)$ \textup{(\textbf{0},~12.1.3.1)} es biyectivo para todo $p$.
\end{corollary}

\begin{proof}
Basta (por \sref[0]{0.12.1.7}) mostrar que los homomorfismos de borde $''E_2^{p0}=H^p(Y,f_*(\mathcal F))\to H^p(X,\mathcal F)$ de la segunda sucesión espectral del funtor compuesto $\Gamma f_*$ son biyectivos.
Pero el término $E_2$ de esta sucesión espectral viene dado por $''E_2^{pq}=H^p(Y,R^qf_*(\mathcal F))$ (G, II, 4.17.1); por el Corolario~\sref{III.1.3.2}, $''E_2^{pq}=0$ para $q>0$, y la sucesión espectral degenera; de ahí nuestra aserción \sref[0]{0.11.1.6}.
\end{proof}

\begin{corollary}[1.3.4]
\label{III.1.3.4}
\label{III.1.3.4-fr}
Sean $f:X\to Y$ un morfismo afín y $g:Y\to Z$ un morfismo.
\oldpage[III]{89}
Para todo $\mathcal O_X$-módulo cuasicoherente $\mathcal F$, el homomorfismo canónico $R^pg_*(f_*(\mathcal F))\to R^p(g\circ f)_*(\mathcal F)$ \textup{(\textbf{0},~12.2.5.1)} es biyectivo para todo $p$.
\end{corollary}

\begin{proof}
Basta observar que, según el Corolario~\sref{III.1.3.3}, para todo abierto afín $W$ de $Z$, el homomorfismo canónico $H^p(g^{-1}(W),f_*(\mathcal F))\to
H^p(f^{-1}(g^{-1}(W)),\mathcal F)$ es biyectivo; esto prueba que el homomorfismo de prehaces que define el homomorfismo canónico $R^pg_*(f_*(\mathcal F))\to R^p(g\circ f)_*(\mathcal F)$ es biyectivo \sref[0]{0.12.2.5}.
\end{proof}

\subsection{Aplicación a la cohomología de preesquemas arbitrarios}
\label{subsection:III.1.4}
\label{subsection:III.1.4-fr}

\begin{proposition}[1.4.1]
\label{III.1.4.1}
\label{III.1.4.1-fr}
Sea $X$ un esquema y sea $\mathfrak U=(U_\alpha)$ un recubrimiento de $X$ por abiertos afines.
Para todo $\mathcal O_X$-módulo cuasicoherente $\mathcal F$, los módulos de cohomología $H^\bullet(X,\mathcal F)$ y $H^\bullet(\mathfrak U,\mathcal F)$ (sobre $\Gamma(X,\mathcal O_X)$) son canónicamente isomorfos.
\end{proposition}

\begin{proof}
Como $X$ es un esquema, toda intersección finita $V$ de abiertos del recubrimiento $\mathfrak U$ es afín \sref[I]{I.5.5.6}; por tanto, $H^q(V,\mathcal F)=0$ para $q\geq1$ por el Teorema~\sref{III.1.3.1}.
La proposición se deduce entonces de un teorema de Leray (G, II, 5.4.1).
\end{proof}

\begin{remark}[1.4.2]
\label{III.1.4.2}
\label{III.1.4.2-fr}
Se observará que el resultado de la Proposición~\sref{III.1.4.1} sigue siendo válido cuando las intersecciones finitas de los conjuntos $U_\alpha$ son afines, incluso si no se supone necesariamente que $X$ sea un esquema.
\end{remark}

\begin{corollary}[1.4.3]
\label{III.1.4.3}
\label{III.1.4.3-fr}
Sea $X$ un esquema de espacio subyacente cuasicompacto, sea $A=\Gamma(X,\mathcal O_X)$ y sea $\mathbf f=(f_i)_{1\leq i\leq r}$ una sucesión finita de elementos de $A$ tal que los $X_{f_i}$ (notación de \sref{III.1.2.1}) sean afines.
Entonces (con las notaciones de \sref{III.1.2.1}), para todo $\mathcal O_X$-módulo cuasicoherente $\mathcal F$, se tiene un isomorfismo canónico funtorial en $\mathcal F$
\begin{equation}
\label{III.1.4.3.1-fr}
\tag{1.4.3.1}
H^q(U,\mathcal F)\xrightarrow{\sim}H^{q+1}((\mathbf f),M)
\qquad\text{para }q\geq1
\end{equation}
y una sucesión exacta funtorial en $\mathcal F$
\begin{equation}
\label{III.1.4.3.2-fr}
\tag{1.4.3.2}
0\longrightarrow H^0((\mathbf f),M)\longrightarrow M
\longrightarrow H^0(U,\mathcal F)
\longrightarrow H^1((\mathbf f),M)\longrightarrow0.
\end{equation}
\end{corollary}

\begin{proof}
Esto se deduce inmediatamente de las Proposiciones~\sref{III.1.4.1} y \sref{III.1.2.3}.
\end{proof}

\begin{env}[1.4.4]
\label{III.1.4.4}
\label{III.1.4.4-fr}
Si $X$ es un \emph{esquema afín}, de la Observación~\sref{III.1.2.5} y de la Proposición~\sref{III.1.4.1} se deduce que, para todo $q\geq0$, los diagramas
\begin{equation}
\label{III.1.4.4.1-fr}
\tag{1.4.4.1}
\xymatrix{
H^q(U,\mathcal F'')\ar[r]^{\partial}\ar[d]
&H^{q+1}(U,\mathcal F')\ar[d]\\
H^{q+1}((\mathbf f),M'')\ar[r]_{\partial}
&H^{q+2}((\mathbf f),M')
}
\end{equation}
correspondientes a una sucesión exacta $0\to\mathcal F'\to\mathcal F\to\mathcal F''\to0$ de $\mathcal O_X$-módulos cuasicoherentes (con las notaciones de la Observación~\sref{III.1.2.5}) son conmutativos.
\end{env}

\begin{proposition}[1.4.5]
\label{III.1.4.5}
\label{III.1.4.5-fr}
\oldpage[III]{90}
Sea $X$ un esquema cuasicompacto, sea $\mathcal L$ un $\mathcal O_X$-módulo invertible, y consideremos el anillo graduado $A_*=\Gamma_*(\mathcal L)$ \sref[0]{0.5.4.6}; entonces $H^\bullet(\mathcal F,\mathcal L)
=\bigoplus_{n\in\mathbb Z}H^\bullet(X,\mathcal F\otimes\mathcal L^{\otimes n})$ es un $A_*$-módulo graduado y, para todo $f\in A_n$, se tiene un isomorfismo canónico
\begin{equation}
\label{III.1.4.5.1-fr}
\tag{1.4.5.1}
H^\bullet(X_f,\mathcal F)
\xrightarrow{\sim}(H^\bullet(\mathcal F,\mathcal L))_{(f)}
\end{equation}
de $(A_*)_{(f)}$-módulos.
\end{proposition}

\begin{proof}
Como $X$ es un esquema cuasicompacto, puede calcularse la cohomología de todos los $\mathcal O_X$-módulos $\mathcal F\otimes\mathcal L^{\otimes n}$ utilizando un mismo recubrimiento finito $\mathfrak U=(U_i)$ por abiertos afines tales que la restricción
$\mathcal L\vert U_i$ es isomorfa a $\mathcal O_X\vert U_i$ para cada
$i$ \sref{III.1.4.1}.
Es entonces inmediato que los $U_i\cap X_f$ son abiertos afines \sref[I]{I.1.3.6}, y puede calcularse así la cohomología $H^\bullet(X_f,\mathcal F\otimes\mathcal L^{\otimes n})$ utilizando el recubrimiento $\mathfrak U\vert X_f=(U_i\cap X_f)$ \sref{III.1.4.1}.
Es inmediato que, para todo $f\in A_n$, la multiplicación por $f$ define un homomorfismo $C^\bullet(\mathfrak U,\mathcal F\otimes\mathcal L^{\otimes m})
\longrightarrow
C^\bullet(\mathfrak U,\mathcal F\otimes\mathcal L^{\otimes(m+n)})$ y, por tanto, un homomorfismo $H^\bullet(\mathfrak U,\mathcal F\otimes\mathcal L^{\otimes m})
\longrightarrow
H^\bullet(\mathfrak U,\mathcal F\otimes\mathcal L^{\otimes(m+n)})$, lo que establece la primera afirmación.
Por otra parte, para un $f\in A_n$ dado, de \sref[I]{I.9.3.2} se deduce que se tiene un isomorfismo de complejos de $(A_*)_{(f)}$-módulos
\[
C^\bullet(\mathfrak U\vert X_f,\mathcal F)
\xrightarrow{\sim}
\left(C^\bullet\left(\mathfrak U,
\bigoplus_{n\in\mathbb Z}\mathcal F\otimes\mathcal L^{\otimes n}
\right)\right)_{(f)}
\]
teniendo en cuenta \sref[I]{I.1.3.9}[ii].
Al pasar a la cohomología de estos complejos se obtiene el isomorfismo (\hyperref[III.1.4.5.1]{1.4.5.1}), recordando que el funtor $M\leadsto M_{(f)}$ es exacto sobre la categoría de $A_*$-módulos graduados.
\end{proof}

\begin{corollary}[1.4.6]
\label{III.1.4.6}
\label{III.1.4.6-fr}
Se supone que se satisfacen las hipótesis de la Proposición~\textup{\sref{III.1.4.5}} y se supone además que $\mathcal L=\mathcal O_X$.
Si se pone $A=\Gamma(X,\mathcal O_X)$, entonces, para todo $f\in A$, se tiene un isomorfismo canónico $H^\bullet(X_f,\mathcal F)
\xrightarrow{\sim}(H^\bullet(X,\mathcal F))_f$ de $A_f$-módulos.
\end{corollary}

\begin{corollary}[1.4.7]
\label{III.1.4.7}
\label{III.1.4.7-fr}
Sea $X$ un esquema cuasicompacto y sea $f$ un elemento de $\Gamma(X,\mathcal O_X)$.
\begin{enumerate}
  \item[{\rm(i)}] Supongamos que el abierto $X_f$ es afín.
    Entonces, para todo $\mathcal O_X$-módulo cuasicoherente $\mathcal F$, todo $i>0$ y todo $\xi\in H^i(X,\mathcal F)$, existe un entero $n>0$ tal que $f^n\xi=0$.
  \item[{\rm(ii)}] Recíprocamente, supongamos que $X_f$ es cuasicompacto y que, para todo haz cuasicoherente de ideales $\mathcal J$ de $\mathcal O_X$ y todo $\zeta\in H^1(X,\mathcal J)$, existe un $n>0$ tal que $f^n\zeta=0$.
    Entonces $X_f$ es afín.
\end{enumerate}
\end{corollary}

\begin{proof}
\medskip\noindent
\begin{enumerate}
  \item[(i)] Si $X_f$ es afín, se tiene $H^i(X_f,\mathcal F)=0$ para todo $i>0$ \sref{III.1.3.1}, de modo que la afirmación se deduce directamente del Corolario~\sref{III.1.4.6}.
  \item[(ii)] En virtud del criterio de Serre \sref[II]{II.5.2.1}, basta probar que, para todo haz cuasicoherente de ideales $\mathcal K$ de $\mathcal O_X\vert X_f$, se tiene $H^1(X_f,\mathcal K)=0$.
    Como $X_f$ es un abierto cuasicompacto de un esquema cuasicompacto $X$, existe un haz cuasicoherente de ideales $\mathcal J$ de $\mathcal O_X$ tal que $\mathcal K=\mathcal J\vert X_f$ \sref[I]{I.9.4.2}.
    Según el Corolario~\sref{III.1.4.6}, se tiene $H^1(X_f,\mathcal K)=(H^1(X,\mathcal J))_f$, y la hipótesis implica que el miembro de la derecha es nulo, lo que da la afirmación.
\end{enumerate}
\end{proof}

\begin{remark}[1.4.8]
\label{III.1.4.8}
\label{III.1.4.8-fr}
Se observará que el Corolario~\sref{III.1.4.7}[i] da una demostración más sencilla de la relación (\textbf{II},~4.5.13.2).
\end{remark}

\begin{lemma}[1.4.9]
\label{III.1.4.9}
\label{III.1.4.9-fr}
\oldpage[III]{91}
Sea $X$ un esquema cuasicompacto, sea $\mathfrak U=(U_i)_{1\leq i\leq n}$ un recubrimiento finito de $X$ por abiertos afines y sea $\mathcal F$ un $\mathcal O_X$-módulo cuasicoherente.
El complejo de haces $\mathcal C^\bullet(\mathfrak U,\mathcal F)$ definido por el recubrimiento $\mathfrak U$ \textup{(G, II, 5.2)} es entonces un $\mathcal O_X$-módulo cuasicoherente.
\end{lemma}

\begin{proof}
De las definiciones (G, II, 5.2) se deduce que $\mathcal C^p(\mathfrak U,\mathcal F)$ es la suma directa de los haces imagen directa de los $\mathcal F\vert U_{i_0\ldots i_p}$ por la inyección canónica $U_{i_0\ldots i_p}\to X$.
La hipótesis de que $X$ es un esquema implica que estas inyecciones son morfismos afines \sref[I]{I.5.5.6}; por tanto, los $\mathcal C^p(\mathfrak U,\mathcal F)$ son cuasicoherentes \sref[II]{II.1.2.6}.
\end{proof}

\begin{proposition}[1.4.10]
\label{III.1.4.10}
\label{III.1.4.10-fr}
Sea $u:X\to Y$ un morfismo separado y cuasicompacto.
Para todo $\mathcal O_X$-módulo cuasicoherente $\mathcal F$, los $R^qu_*(\mathcal F)$ son $\mathcal O_Y$-módulos cuasicoherentes.
\end{proposition}

\begin{proof}
La cuestión es local sobre $Y$, de modo que puede suponerse que $Y$ es afín.
Entonces $X$ es una unión finita de abiertos afines $U_i$ ($1\leq i\leq n$); sea $\mathfrak U$ el recubrimiento $(U_i)$.
Además, como $Y$ es un esquema, de \sref[I]{I.5.5.10} se deduce que, para todo abierto afín $V\subset Y$, la inyección canónica $u^{-1}(V)\to X$ es un morfismo afín; se concluye (Proposición~\sref{III.1.4.1} y (G, II, 5.2)) que se tiene un isomorfismo canónico
\begin{equation}
\label{III.1.4.10.1-fr}
\tag{1.4.10.1}
H^\bullet(u^{-1}(V),\mathcal F)
\xrightarrow{\sim}H^\bullet(\Gamma(V,\mathcal K^\bullet))
\end{equation}
donde se pone $\mathcal K^\bullet=u_*(\mathcal C^\bullet(\mathfrak U,\mathcal F))$.
Según el Lema~\sref{III.1.4.9} y \sref[I]{I.9.2.2}, $\mathcal K^\bullet$ es un $\mathcal O_Y$-módulo cuasicoherente; además, constituye un \emph{complejo de haces}, puesto que también lo es $\mathcal C^\bullet(\mathfrak U,\mathcal F)$.
De la definición de la cohomología $\mathcal H^\bullet(\mathcal K^\bullet)$ (G, II, 4.1) se deduce entonces que esta última está formada por $\mathcal O_Y$-módulos cuasicoherentes \sref[I]{I.4.1.1}.
Como (para $V$ afín en $Y$) el funtor $\Gamma(V,\mathcal G)$ es exacto en $\mathcal G$ sobre la categoría de $\mathcal O_Y$-módulos cuasicoherentes, se tiene (G, II, 4.1)
\begin{equation}
\label{III.1.4.10.2-fr}
\tag{1.4.10.2}
H^\bullet(\Gamma(V,\mathcal K^\bullet))
=\Gamma(V,\mathcal H^\bullet(\mathcal K^\bullet)).
\end{equation}
Observemos finalmente que resulta de la definición del homomorfismo canónico \[
H^\bullet(\mathfrak U,\mathcal F)\longrightarrow H^\bullet(X,\mathcal F)
\]
dado en (G, II, 5.2), se deduce que, si $V'\subset V$ es otro abierto afín de $Y$, el
diagrama
\[
\xymatrix{
H^\bullet(u^{-1}(V),\mathcal F)\ar[r]^{\sim}\ar[d]
&H^\bullet(\Gamma(V,\mathcal K^\bullet))\ar[d]\\
H^\bullet(u^{-1}(V'),\mathcal F)\ar[r]^{\sim}
&H^\bullet(\Gamma(V',\mathcal K^\bullet))
}
\]
es conmutativo.
De lo anterior se concluye que los isomorfismos (\hyperref[III.1.4.10.1]{1.4.10.1}) definen un isomorfismo de $\mathcal O_Y$-módulos
\begin{equation}
\label{III.1.4.10.3-fr}
\tag{1.4.10.3}
R^\bullet u_*(\mathcal F)
\xrightarrow{\sim}\mathcal H^\bullet(\mathcal K^\bullet)
\end{equation}
\oldpage[III]{92}
y, por consiguiente, $R^\bullet u_*(\mathcal F)$ es cuasicoherente.
\end{proof}

Además, de (\hyperref[III.1.4.10.3]{1.4.10.3}), (\hyperref[III.1.4.10.2]{1.4.10.2}) y (\hyperref[III.1.4.10.1]{1.4.10.1}) se deduce que:
\begin{corollary}[1.4.11]
\label{III.1.4.11}
\label{III.1.4.11-fr}
Bajo las hipótesis de la Proposición~\sref{III.1.4.10}, para todo abierto afín $V$ de $Y$, el homomorfismo canónico
\begin{equation}
\label{III.1.4.11.1-fr}
\tag{1.4.11.1}
H^q(u^{-1}(V),\mathcal F)\longrightarrow
\Gamma(V,R^qu_*(\mathcal F))
\end{equation}
es un isomorfismo para todo $q\geq0$.
\end{corollary}

\begin{corollary}[1.4.12]
\label{III.1.4.12-fr}
Supongamos verificadas las hipótesis de \textup{(\hyperref[III.1.4.10-fr]{1.4.10})} y supongamos además que $Y$ es cuasicompacto.
Entonces existe un entero $r>0$ tal que, para todo $\mathcal O_X$-módulo cuasicoherente $\mathcal F$ y todo entero $q>r$, se tiene $R^qu_*(\mathcal F)=0$.
Si $Y$ es afín, puede tomarse como $r$ un entero tal que exista un recubrimiento de $X$ formado por $r$ abiertos afines.
\end{corollary}

\begin{proof}
Como puede recubrirse $Y$ por un número finito de abiertos afines, en virtud del Corolario~\sref{III.1.4.11} puede reducirse la demostración a la segunda afirmación.
Si $\mathfrak U$ es un recubrimiento de $X$ por $r$ abiertos afines, se tiene $H^q(\mathfrak U,\mathcal F)=0$ para $q>r$, puesto que las cocadenas de $C^q(\mathfrak U,\mathcal F)$ son alternadas; la afirmación se deduce entonces de la Proposición~\sref{III.1.4.1}.
\end{proof}

\begin{corollary}[1.4.13]
\label{III.1.4.13-fr}
Supongamos verificadas las hipótesis de \textup{(\hyperref[III.1.4.10-fr]{1.4.10})} y supongamos además que $Y=\operatorname{Spec}(A)$ es afín.
Entonces, para todo $\mathcal O_X$-módulo cuasicoherente $\mathcal F$ y todo $f\in A$, se tiene
\[
\Gamma(Y_f,R^qu_*(\mathcal F))
=(\Gamma(Y,R^qu_*(\mathcal F)))_f
\]
salvo isomorfismo canónico.
\end{corollary}

\begin{proof}
Esto se deduce del hecho de que $R^qu_*(\mathcal F)$ es un $\mathcal O_Y$-módulo cuasicoherente \sref[I]{I.1.3.7}.
\end{proof}

\begin{proposition}[1.4.14]
\label{III.1.4.14}
\label{III.1.4.14-fr}
Sean $f:X\to Y$ un morfismo separado y cuasicompacto y $g:Y\to Z$ un morfismo afín.
Para todo $\mathcal O_X$-módulo cuasicoherente $\mathcal F$, el homomorfismo canónico $R^p(g\circ f)_*(\mathcal F)\to g_*(R^pf_*(\mathcal F))$ \textup{(\textbf{0},~12.2.5.2)} es biyectivo para todo $p$.
\end{proposition}

\begin{proof}
Para todo abierto afín $W$ de $Z$, $g^{-1}(W)$ es un abierto afín de $Y$.
El homomorfismo de prehaces que define el homomorfismo canónico
\[
R^p(g\circ f)_*(\mathcal F)\longrightarrow g_*(R^pf_*(\mathcal F))
\]
\sref[0]{0.12.2.5} es, por tanto, biyectivo en virtud del Corolario~\sref{III.1.4.11}.
\end{proof}

\begin{proposition}[1.4.15]
\label{III.1.4.15}
\label{III.1.4.15-fr}
Sean $u:X\to Y$ un morfismo separado de tipo finito y $v:Y'\to Y$ un morfismo plano de preesquemas \sref[0]{0.6.7.1}; sea $u'=u_{(Y')}$, de modo que se tiene el diagrama conmutativo
\begin{equation}
\label{III.1.4.15.1-fr}
\tag{1.4.15.1}
\xymatrix{
X\ar[d]_u & X'=X_{(Y')}\ar[l]_{v'}\ar[d]^{u'}\\
Y & Y'\ar[l]^v
}\quad.
\end{equation}
Entonces, para todo $\mathcal O_X$-módulo cuasicoherente $\mathcal F$, $R^qu_*'(\mathcal F')$, donde
$\mathcal F'=v'^*(\mathcal F)=\mathcal F\otimes_{\mathcal O_Y}\mathcal O_{Y'}$,
es canónicamente isomorfo a
$R^qu_*(\mathcal F)\otimes_{\mathcal O_Y}\mathcal O_{Y'}
=v^*(R^qu_*(\mathcal F))$ para todo $q\geq0$.
\end{proposition}

\oldpage[III]{93}
\begin{proof}
El homomorfismo canónico
$\rho:\mathcal F\to v_*'(v'^*(\mathcal F))$ (\textbf{0},~4.4.3.2) define por funtorialidad un homomorfismo
\begin{equation}
\label{III.1.4.15.2-fr}
\tag{1.4.15.2}
R^qu_*(\mathcal F)\longrightarrow R^qu_*(v_*'(\mathcal F')).
\end{equation}

Por otra parte, poniendo $w=u\circ v'=v\circ u'$, se tienen los homomorfismos canónicos (\textbf{0}~,12.2.5.1 y 12.2.5.2)
\begin{equation}
\label{III.1.4.15.3-fr}
\tag{1.4.15.3}
R^qu_*(v_*'(\mathcal F'))\longrightarrow R^qw_*(\mathcal F')
\longrightarrow v_*(R^qu_*'(\mathcal F')).
\end{equation}

Componiendo (\hyperref[III.1.4.15.3]{1.4.15.3}) y (\hyperref[III.1.4.15.2]{1.4.15.2}), se obtiene un homomorfismo
\[
\psi:R^qu_*(\mathcal F)\longrightarrow v_*(R^qu_*'(\mathcal F'))
\]
y se obtiene finalmente un homomorfismo canónico (cuya definición no impone \emph{hipótesis alguna} sobre $v$)
\begin{equation}
\label{III.1.4.15.4-fr}
\tag{1.4.15.4}
\psi^\sharp:v^*(R^qu_*(\mathcal F))\longrightarrow R^qu_*'(\mathcal F')
\end{equation}
y hay que probar que es un isomorfismo cuando $v$ es \emph{plano}.
Es claro que la cuestión es local sobre $Y$ y $Y'$; puede suponerse, por tanto, que $Y=\operatorname{Spec}(A)$ y $Y'=\operatorname{Spec}(B)$; utilizaremos además el
\begin{lemma}[1.4.15.5]
\label{III.1.4.15.5-fr}
Sean $\varphi:A\to B$ un homomorfismo de anillos, $Y=\operatorname{Spec}(A)$, $X=\operatorname{Spec}(B)$, $f:X\to Y$ el morfismo correspondiente a $\varphi$, y $M$ un $B$-módulo.
Para que el $\mathcal O_X$-módulo $\widetilde M$ sea $f$-plano \sref[0]{0.6.7.1}, es necesario y suficiente que $M$ sea un $A$-módulo plano.
En particular, para que el morfismo $f$ sea plano, es necesario y suficiente que $B$ sea un $A$-módulo plano.
\end{lemma}

Esto se deduce de la definición \sref[0]{0.6.7.1} y de \sref[0]{0.6.3.3}, teniendo en cuenta \sref[I]{I.1.3.4}.

Así pues, de (\hyperref[III.1.4.11.1]{1.4.11.1}) y de las definiciones de los homomorfismos (\hyperref[III.1.4.15.3]{1.4.15.3}) (cf.~\sref[0]{0.12.2.5}) se deduce que $\psi$ corresponde entonces al morfismo compuesto
\[
H^q(X,\mathcal F)\xrightarrow{\rho_q}
H^q(X,v_*'(v'^*(\mathcal F)))\xrightarrow{\theta_q}
H^q(X',v'^*(v_*'(v'^*(\mathcal F))))\xrightarrow{\sigma_q}
H^q(X',v'^*(\mathcal F))
\]
donde $\rho_q$ y $\sigma_q$ son los homomorfismos de cohomología correspondientes a los morfismos canónicos $\rho$ y $\sigma:v'^*(v_*'(\mathcal G'))\to\mathcal G'$, y $\theta_q$ es el $\varphi$-morfismo (\textbf{0},~12.1.3.1) relativo al $\mathcal O_X$-módulo $v_*'(v'^*(\mathcal F))$.
Pero, por la funtorialidad de $\theta_q$, se tiene el diagrama conmutativo
\[
\xymatrix{
H^q(X,\mathcal F)\ar[rr]^{\rho_q}\ar[dd]_{\theta_q}
&&H^q(X,v_*'(v'^*(\mathcal F)))\ar[dd]^{\theta_q}\\\\
H^q(X',v'^*(\mathcal F))\ar[rr]_{v'^*(\rho_q)}
&&H^q(X',v'^*(v_*'(v'^*(\mathcal F))))
}
\]
\oldpage[III]{94}
y, como por definición \sref[0]{0.4.4.3} $v'^*(\rho)$ es el inverso de $\sigma$, se ve que el morfismo compuesto considerado anteriormente no es finalmente otro que $\theta_q$; por consiguiente, $\psi^\sharp$ es el $B$-homomorfismo asociado $H^q(X,\mathcal F)\otimes_A B\to H^q(X',\mathcal F')$.
Como $u$ es de tipo finito, $X$ es una unión finita de abiertos afines $U_i$ ($1\leq i\leq r$); sea $\mathfrak U$ el recubrimiento $(U_i)$.
Como $v$ es un morfismo afín, también lo es $v'$ \sref[II]{II.1.6.2}[iii]; por consiguiente, los $U_i'=v'^{-1}(U_i)$ forman un recubrimiento abierto afín $\mathfrak U'$ de $X'$.
Se sabe entonces \sref[0]{0.12.1.4.2} que el diagrama
\[
\xymatrix{
H^q(\mathfrak U,\mathcal F)\ar[r]^{\theta_q}\ar[d]
&H^q(\mathfrak U',\mathcal F')\ar[d]\\
H^q(X,\mathcal F)\ar[r]_{\theta_q}&H^q(X',\mathcal F')
}
\]
es conmutativo, y las flechas verticales son isomorfismos puesto que $X$ y $X'$ son esquemas \sref{III.1.4.1}.
Por consiguiente, basta probar que el $\varphi$-morfismo canónico $\theta_q:H^q(\mathfrak U,\mathcal F)\to H^q(\mathfrak U',\mathcal F')$ es tal que el $B$-homomorfismo asociado
\[
H^q(\mathfrak U,\mathcal F)\otimes_A B
\longrightarrow H^q(\mathfrak U',\mathcal F')
\]
sea un isomorfismo.
Para toda sucesión $\mathbf s=(i_k)_{0\leq k\leq p}$ de $p+1$ índices de $[1,r]$, pongamos $U_{\mathbf s}=\bigcap\limits_{k=0}^pU_{i_k}$, $U_{\mathbf s}'=\bigcap\limits_{k=0}^{\boldsymbol{p}}U_{i_k}'=v'^{-1}(U_{\mathbf s})$, $M_{\mathbf s}=\Gamma(U_{\mathbf s},\mathcal F)$ y $M_{\mathbf s}'=\Gamma(U_{\mathbf s}',\mathcal F')$.
La aplicación canónica $M_{\mathbf s}\otimes_A B\to M_{\mathbf s}'$ es un isomorfismo \sref[I]{I.1.6.5}; por tanto, la aplicación canónica $C^p(\mathfrak U,\mathcal F)\otimes_A B\to
C^p(\mathfrak U',\mathcal F')$ es un isomorfismo, mediante el cual $d\otimes1$ se identifica con la aplicación coborde $C^p(\mathfrak U',\mathcal F')\to C^{p+1}(\mathfrak U',\mathcal F')$.
Como $B$ es un $A$-módulo \emph{plano}, de la definición de los módulos de cohomología se deduce que la aplicación canónica $H^q(\mathfrak U,\mathcal F)\otimes_A B\to
H^q(\mathfrak U',\mathcal F')$ es un isomorfismo \sref[0]{0.6.1.1}.
Este resultado se generalizará más adelante en \EGAUnavailableHyperref{section:III.6}{\textsection6}.
\end{proof}

\begin{corollary}[1.4.16]
\label{III.1.4.16}
\label{III.1.4.16-fr}
Sean $A$ un anillo, $X$ un $A$-esquema de tipo finito y $B$ una $A$-álgebra fielmente plana sobre $A$.
Para que $X$ sea afín, es necesario y suficiente que $X\otimes_A B$ lo sea.
\end{corollary}

\begin{proof}
La condición es evidentemente necesaria \sref[I]{I.3.2.2}; mostremos que es suficiente.
Como $X$ es separado sobre $A$ y el morfismo $\operatorname{Spec}(B)\to\operatorname{Spec}(A)$ es plano, de la Proposición~\sref{III.1.4.1} se deduce que
\begin{equation}
\label{III.1.4.16.1-fr}
\tag{1.4.16.1}
H^i(X\otimes_A B,\mathcal F\otimes_A B)
=H^i(X,\mathcal F)\otimes_A B
\end{equation}
para todo $i\geq0$ y todo $\mathcal O_X$-módulo cuasicoherente $\mathcal F$.
Si $X\otimes_A B$ es afín, el miembro izquierdo de (\hyperref[III.1.4.16.1]{1.4.16.1}) es nulo para $i=1$; por tanto, también lo es $H^1(X,\mathcal F)$, puesto que $B$ es un $A$-módulo fielmente plano.
Como $X$ es un esquema cuasicompacto, la demostración termina por el criterio de Serre \sref[II]{II.5.2.1}.
\end{proof}

\begin{proposition}[1.4.17]
\label{III.1.4.17}
\label{III.1.4.17-fr}
Sea $X$ un preesquema y sea $0\to\mathcal F\xrightarrow{u}\mathcal G\xrightarrow{v}\mathcal H\to0$ una sucesión exacta de $\mathcal O_X$-módulos.
Si $\mathcal F$ y $\mathcal H$ son cuasicoherentes, también lo es $\mathcal G$.
\end{proposition}

\oldpage[III]{95}
\begin{proof}
La cuestión es local sobre $X$, de modo que puede suponerse que $X=\operatorname{Spec}(A)$ es afín; basta entonces probar que $\mathcal G$ satisface las condiciones (\textit{d}1) y (\textit{d}2) de \sref[I]{I.1.4.1} (con $V=X$).
La verificación de (\textit{d}2) es inmediata, pues, si $t\in\Gamma(X,\mathcal G)$ se anula al restringirlo a $D(f)$, lo mismo ocurre con su imagen $v(t)\in\Gamma(X,\mathcal H)$; por tanto, existe un $m>0$ tal que $f^mv(t)=v(f^mt)=0$ \sref[I]{I.1.4.1}, y, como $\Gamma$ es exacto por la izquierda, $f^mt=u(s)$, donde $s\in\Gamma(X,\mathcal F)$; como $u$ es inyectivo, la restricción de $s$ a $D(f)$ es nula, luego por \sref[I]{I.1.4.1} existe un entero $n>0$ tal que $f^ns=0$; se deduce finalmente que $f^{m+n}t=u(f^ns)=0$.

Compruébese ahora (\textit{d}2); sea $t'\in\Gamma(D(f),\mathcal G)$; como $\mathcal H$ es cuasicoherente, existe un entero $m$ tal que $f^mv(t')=v(f^mt')$ se prolonga a una sección $z\in\Gamma(X,\mathcal H)$ \sref[I]{I.1.4.1}.
Pero, en virtud del Teorema~\sref{III.1.3.1} (o de (\textbf{I},~5.1.9.2)) aplicado al $\mathcal O_X$-módulo cuasicoherente $\mathcal F$, la sucesión $\Gamma(X,\mathcal G)\to\Gamma(X,\mathcal H)\to0$ es exacta; existe, pues, $t\in\Gamma(X,\mathcal G)$ tal que $z=v(t)$; se ve así que $v(f^mt'\setmin t'')=0$, denotando por $t''$ la restricción de $t$ a $D(f)$; por tanto, $f^mt'\setmin t''=u(s')$, donde $s'\in\Gamma(D(f),\mathcal F)$.
Pero, como $\mathcal F$ es cuasicoherente, existe un entero $n>0$ tal que $f^ns'$ se prolonga a una sección $s\in\Gamma(X,\mathcal F)$; como $f^{m+n}t'\setmin f^nt''=u(f^ns')$, se ve que $f^{m+n}t'$ es la restricción a $D(f)$ de una sección $f^nt+u(f^ns)\in\Gamma(X,\mathcal G)$, lo que termina la demostración.
\end{proof}
